% Palette Introduction --- Chinese.
% Auto-generated by gen_palette_intro.py.  Do not edit by hand.
\chapter{调色板介绍}\label{ch:palette-intro}

本章按编号依次介绍 24 组主题调色板。每一组都保留三段固定结构——\emph{色彩哲学}阐明配色背后的世界观,\emph{六色配色组}列出 \texttt{pal1} 到 \texttt{pal6} 的色值与释义,随后的小节(\emph{设计思路}、\emph{色彩结构}或\emph{美学法则}等)再对每一色或每一条结构性原则展开说明。配色按编号顺序排列、自然分页,不再强制每页限定数量,只保证六色配色表保持完整、不被分割。

% ---------------------------------------------------------------------------
\bookpalettecardhead{0\enspace OUC Default}
\palettesubsec{色彩哲学}
\noindent 这是模板出厂时的底色。蓝是青年学者的颜色——从清晨第一页讲义、午后图书馆窗台,到深夜实验台的冷光,蓝构成学术工作的整个色温。六层由浅入深的蓝刻意维持理性的克制,模拟一本书从翻开到合上的全过程;它不是一组要被注视的色,而是一组让文字自身被看见的色。这是写作者与读者之间最不打扰的契约,也是中国海洋大学海德学院 LaTeX 模板希望被使用者长期相处的方式。

\begin{fullwidth}\small\noindent
\begin{tabularx}{\linewidth}{@{}clllX@{}}
\toprule\rowcolor{booktableheadcolor}
\booktableheadcell{} & \booktableheadcell{中文名} & \booktableheadcell{原名} & \booktableheadcell{RGB} & \booktableheadcell{释义} \\
\midrule
\bookcardswatch{239,246,255} & 纸海蓝 & Paper Foam & 239, 246, 255 & 翻开扉页时纸面反射的最浅一抹蓝,几近白 \\
\bookcardswatch{219,234,254} & 冰晶蓝 & Desert Foam & 219, 234, 254 & 实验室冷光下纸张投出的浅影,清洁而无负担 \\
\bookcardswatch{96,165,250} & 天玻璃蓝 & Sea Smoke & 96, 165, 250 & 窗外的天,是学者抬头喘息时的呼吸色 \\
\bookcardswatch{37,99,235} & 讲堂蓝 & Crisp Sky & 37, 99, 235 & 黑板上一笔粉笔勾出的强调,是论证之色 \\
\bookcardswatch{30,58,138} & 深思蓝 & Twilight Blue & 30, 58, 138 & 思考最深处的暗蓝,接近梦的边界 \\
\bookcardswatch{30,41,59} & 铅墨石 & Midnight Hull & 30, 41, 59 & 钢笔留在笔记本上的最后一行字 \\
\bottomrule
\end{tabularx}\end{fullwidth}

\palettesubsec{阅读弧线的色彩映射}
\begin{palettebullets}
\item \textbf{Paper Foam} --- 纸海蓝是「开卷」的呼吸口,几近于白,允许任何一种文字以最少的视觉负担进入读者的眼睛
\item \textbf{Frost Crystal} --- 冰晶蓝承担背景的次轻调,在长时间阅读中替代纯白,降低疲劳而不引起注意
\item \textbf{Sky Glass} --- 天玻璃蓝是六色中唯一带有「自然」气息的中调,把书页从全然的人工色阵列里拉回一次窗外的视野
\item \textbf{Lecture Sapphire} --- 讲堂蓝承担论证的强调职责,是六色中唯一允许「响」的色——但响得克制,像粉笔在板上敲下重音
\item \textbf{Reverie Navy} --- 深思蓝把思考压向最深处,适合做长引文、关键定理、章节扉页的暗调主色
\item \textbf{Slate Ink} --- 铅墨石几乎是黑,却保有蓝的内核,是正文文字与所有结构线条最终落笔的颜色——书在它的颜色里收束
\end{palettebullets}
\par\bigskip
% ---------------------------------------------------------------------------
\bookpalettecardhead{1\enspace Brunneophobia}
\palettesubsec{色彩哲学}
\noindent 名字直译为「恐褐症」,却用整组褐色直面这种偏见。褐不是廉价、不是过时,而是被现代审美刻意冷落的暖。从沙土、橘黄、铁锈到深焦,六色一路下沉,如同走进一间堆满旧书的木质工作室,鼻腔里有皮革、烟草、湿木的气息。色彩的勇气不在于鲜艳,而在于敢做被回避的那一种;这组配色把所有被嫌弃的褐重新收编进一种成熟的体面。

\begin{fullwidth}\small\noindent
\begin{tabularx}{\linewidth}{@{}clllX@{}}
\toprule\rowcolor{booktableheadcolor}
\booktableheadcell{} & \booktableheadcell{中文名} & \booktableheadcell{原名} & \booktableheadcell{RGB} & \booktableheadcell{释义} \\
\midrule
\bookcardswatch{238,211,180} & 沙浪米 & Lighthouse & 238, 211, 180 & 沙漠日出的最浅暖白,是褐色系的「呼吸口」 \\
\bookcardswatch{213,148,79} & 烟橘 & Sand Lily & 213, 148, 79 & 烟雾掩映下的橘黄,温暖却不张扬 \\
\bookcardswatch{213,148,79} & 蜂蜡金 & Reef Coral & 213, 148, 79 & 同色重复——表明这组褐刻意收紧色相 \\
\bookcardswatch{180,69,15} & 铁锈赤 & Tidal Mauve & 180, 69, 15 & 暴露在湿气中数十年的铁,带着血与土 \\
\bookcardswatch{86,67,53} & 旧皮褐 & Inkwell & 86, 67, 53 & 老沙发与皮装订书脊的颜色,有岁月气息 \\
\bookcardswatch{42,23,14} & 焦土黑 & Deep Tide & 42, 23, 14 & 篝火熄灭后炭灰里残留的最深褐 \\
\bottomrule
\end{tabularx}\end{fullwidth}

\palettesubsec{褐色家族的内部结构}
\begin{palettebullets}
\item \textbf{Desert Foam} --- 沙浪米作为整组配色的最浅调,承担所有「呼吸」与「留白」的职责,是褐色系里唯一不发声的色
\item \textbf{Smoke Tangerine} --- 烟橘负责中段的温度,把整组配色从严肃拉回日常,是这组色与读者之间最亲切的接触面
\item \textbf{Beeswax Honey} --- 蜂蜡金与烟橘共用色相,是整组配色刻意为之的「收紧」——意在证明克制本身即是力量
\item \textbf{Rust Crimson} --- 铁锈赤是这组色彩的「锋」——它把柔和的褐推向有侵略性的红,是唯一允许刺痛的色
\item \textbf{Aged Leather} --- 旧皮褐承担文献质感,是图书馆、档案、皮装订所共享的深调,把读物的物质感留在了色彩里
\item \textbf{Scorched Earth} --- 焦土黑是结构的根,所有上层暖色都因为它的存在而得以稳稳站立——克服恐褐症的最终一步,是承认黑的褐属性
\end{palettebullets}
\par\bigskip
% ---------------------------------------------------------------------------
\bookpalettecardhead{2\enspace Van Dyke}
\palettesubsec{色彩哲学}
\noindent 致敬十七世纪佛兰德画家 Anthony van Dyck——他的肖像画用「凡戴克褐」打底,以薄釉层层堆叠出贵族肌肤的透光感。这组配色从他的画布生长而来:粉调的肌肤、淡紫的丝绸、暗红的天鹅绒、深褐的背景,皆是巴洛克肖像的标准物质。六色之间没有戏剧冲突,只有褪色后的体面;它们组成的不是一幅画,而是一幅画被时间洗过之后,墙上仍然挂着的那种气场。

\begin{fullwidth}\small\noindent
\begin{tabularx}{\linewidth}{@{}clllX@{}}
\toprule\rowcolor{booktableheadcolor}
\booktableheadcell{} & \booktableheadcell{中文名} & \booktableheadcell{原名} & \booktableheadcell{RGB} & \booktableheadcell{释义} \\
\midrule
\bookcardswatch{236,194,188} & 肌肤粉 & Vleeskleur & 236, 194, 188 & 巴洛克肖像中贵族面颊的薄釉粉调 \\
\bookcardswatch{169,159,191} & 丝绸藤 & Zijden Sering & 169, 159, 191 & 折光丝绸在阴影里的淡紫,优雅而冷 \\
\bookcardswatch{169,159,191} & 烟雾紫 & Rooklila & 169, 159, 191 & 同前——丝绸两面在画布上几乎不分 \\
\bookcardswatch{191,113,133} & 玫瑰呢 & Rozenfluweel & 191, 113, 133 & 天鹅绒礼服里渗出的玫瑰红,沉而不艳 \\
\bookcardswatch{68,60,94} & 暗紫绒 & Pruimenpurper & 68, 60, 94 & 背景帷幕的深紫,把光衬得更亮 \\
\bookcardswatch{61,43,39} & 凡戴克褐 & Van Dijck-bruin & 61, 43, 39 & 该画家的标志底色,以沥青颜料调和而成 \\
\bottomrule
\end{tabularx}\end{fullwidth}

\palettesubsec{肖像深度的色彩映射}
\begin{palettebullets}
\item \textbf{Carnation Skin} --- 肌肤粉是画面最前景——画中人物的面颊、手指、颈侧,是观看者第一眼接触到的「人」
\item \textbf{Silken Wisteria} --- 丝绸藤是衣料的高光,把柔软的折光感留在了色相中,优雅而带着距离
\item \textbf{Smoke Lilac} --- 烟雾紫与丝绸藤同色,刻意制造重复,呼应巴洛克绘画里对单色叠层的偏爱
\item \textbf{Rose Velvet} --- 玫瑰呢是中景礼服的颜色,把人物从画中抬起、推近,是这组配色的色温支点
\item \textbf{Plum Plush} --- 暗紫绒是背景帷幕,负责把前景的人物从黑暗中托出,自己几乎不被注意
\item \textbf{Van Dyke Brown} --- 凡戴克褐是整幅画的「地」——是颜料、画布、画框、墙面共同沉淀下来的底色,任何高光都因它而立
\end{palettebullets}
\par\bigskip
% ---------------------------------------------------------------------------
\bookpalettecardhead{3\enspace Back in Black}
\palettesubsec{色彩哲学}
\noindent 致敬 AC/DC 1980 年的同名专辑,但配色本身比摇滚更克制。粉与灰柔化了黑的攻击性,使整组色更接近后朋克美学:皮夹克、烟雾、舞台灯熄灭后那一刻的余温。粉是疤痕,黑是夜——它们之间所有的灰,是青春过后还留在身体里的回声。这是一组「成熟版的叛逆」,把锋芒收进了体面之内。

\begin{fullwidth}\small\noindent
\begin{tabularx}{\linewidth}{@{}clllX@{}}
\toprule\rowcolor{booktableheadcolor}
\booktableheadcell{} & \booktableheadcell{中文名} & \booktableheadcell{原名} & \booktableheadcell{RGB} & \booktableheadcell{释义} \\
\midrule
\bookcardswatch{240,217,228} & 婴儿粉 & Powder Pink & 240, 217, 228 & 黑暗里唯一不设防的颜色,是叛逆的反面 \\
\bookcardswatch{193,160,172} & 灰玫瑰 & Dust Rose & 193, 160, 172 & 蒙了一层烟的粉,失去鲜艳后的体面 \\
\bookcardswatch{193,160,172} & 烬粉 & Ember Pink & 193, 160, 172 & 重复以加强中调,粉的两种语义并置 \\
\bookcardswatch{128,108,121} & 雾紫灰 & Smoky Mauve & 128, 108, 121 & 舞台灯下烟雾的灰紫,既不明也不暗 \\
\bookcardswatch{74,63,75} & 黑天鹅绒 & Black Velvet & 74, 63, 75 & 皮夹克的暗影,介于黑与紫之间 \\
\bookcardswatch{22,19,21} & 夜墨 & Inkwell Black & 22, 19, 21 & 一组配色最沉的锚点,几乎黑 \\
\bottomrule
\end{tabularx}\end{fullwidth}

\palettesubsec{后朋克美学的色彩映射}
\begin{palettebullets}
\item \textbf{Powder Pink} --- 婴儿粉是这组配色的「不设防」——它是黑底之上唯一一个温柔的入口,叛逆的反面才是叛逆的内核
\item \textbf{Dust Rose} --- 灰玫瑰把粉的鲜艳压了下去,是青春之后留下的余韵,优雅地承认了「不再年轻」
\item \textbf{Ember Pink} --- 烬粉与灰玫瑰共享色相,是「重复即强调」的视觉修辞,把中调的灰玫瑰提升为整组色的真正核心
\item \textbf{Smoky Mauve} --- 雾紫灰是舞台烟雾的颜色,过渡了粉与黑,使两端的对立看起来像同一种情绪的两面
\item \textbf{Black Velvet} --- 黑天鹅绒是皮夹克的暗影,带着紫的呼吸——黑不必是死的,它可以是一种被穿在身上的物质
\item \textbf{Inkwell Black} --- 夜墨是最沉的锚点,接近纯黑却保有微紫底调,是整组色的句号——但句号之后还可以有下一首歌
\end{palettebullets}
\par\bigskip
% ---------------------------------------------------------------------------
\bookpalettecardhead{4\enspace Belle of the Ball}
\palettesubsec{色彩哲学}
\noindent 舞会上那位最被注视的女子从不穿原色——她穿的是被无数烛光暖化过的旧色:陶土、铜锈、橄榄、绛红。这组配色拒绝二十世纪的合成染料,从十九世纪欧洲晚宴厅与殖民地热带气候的混合美学中提取——壁炉、桃花心木、暮色窗帘、染了酒渍的桌布。被注视并不靠鲜艳,而靠所有颜色都已被时间柔化过的那种「被生活过的高级感」。

\begin{fullwidth}\small\noindent
\begin{tabularx}{\linewidth}{@{}clllX@{}}
\toprule\rowcolor{booktableheadcolor}
\booktableheadcell{} & \booktableheadcell{中文名} & \booktableheadcell{原名} & \booktableheadcell{RGB} & \booktableheadcell{释义} \\
\midrule
\bookcardswatch{226,203,192} & 香槟珠光 & Champagne Pearl & 226, 203, 192 & 舞会大厅烛光下女子肌肤的珠光底色 \\
\bookcardswatch{206,171,150} & 玫瑰陶土 & Rose Terracotta & 206, 171, 150 & 复古口红与陶土色的中间调,被注视的暖色 \\
\bookcardswatch{210,135,106} & 陶釉橘 & Glazed Clay & 210, 135, 106 & 陶器釉面的橘红,温润而成熟 \\
\bookcardswatch{229,74,57} & 辣椒朱 & Chilli Scarlet & 229, 74, 57 & 热带气候里晾在阳光下的辣椒色,是这组配色的「亮点」 \\
\bookcardswatch{118,118,44} & 橄榄绿 & Olive Verdigris & 118, 118, 44 & 橄榄叶被氧化后的中性绿,带些铜锈 \\
\bookcardswatch{53,77,4} & 苔藓深绿 & Moss Deepwood & 53, 77, 4 & 桃花心木家具与深色窗帘的森林暗调 \\
\bottomrule
\end{tabularx}\end{fullwidth}

\palettesubsec{被注视者的色彩策略}
\begin{palettebullets}
\item \textbf{Champagne Pearl} --- 香槟珠光是被烛光暖化的肌肤底色,是整组配色的「呼吸」,也是被注视者最不设防的一层
\item \textbf{Rose Terracotta} --- 玫瑰陶土承担「复古性感」,是被时间筛过的口红色,既是装饰也是身份的证明
\item \textbf{Glazed Clay} --- 陶釉橘是手作之色,把工业感从配色中排除,留下手心温度与陶土的物质感
\item \textbf{Chilli Scarlet} --- 辣椒朱是全组唯一允许「亮」的色,是裙摆的红、戒指的红、唇上最后一笔的红——但只能小面积出现
\item \textbf{Olive Verdigris} --- 橄榄绿是这组色彩的冷调支点,把所有热带暖色拉回欧洲晚宴厅的克制,是配色的色温调节器
\item \textbf{Moss Deepwood} --- 苔藓深绿是背景与家具的暗调,是这组配色的「重」,没有它,所有的亮都会失去支点
\end{palettebullets}
\par\bigskip
% ---------------------------------------------------------------------------
\bookpalettecardhead{5\enspace Pine Tree}
\palettesubsec{色彩哲学}
\noindent 以一棵北方松树作为坐标,从树冠到树根逐层收色:阳光打在新生针叶上的暖金、被风吹松动的旧针、午后树荫里的玫瑰陶土、树干裂痕里渗出的赭色,以及最底层埋在腐殖土里的森林暗墨。配色逻辑刻意离开「绿色」——一棵真正的松树在视觉上远比绿色复杂得多,它的颜色其实是阳光、岁月、土壤、伤口与影子共同写下的一段年表。

\begin{fullwidth}\small\noindent
\begin{tabularx}{\linewidth}{@{}clllX@{}}
\toprule\rowcolor{booktableheadcolor}
\booktableheadcell{} & \booktableheadcell{中文名} & \booktableheadcell{原名} & \booktableheadcell{RGB} & \booktableheadcell{释义} \\
\midrule
\bookcardswatch{238,200,111} & 阳针金 & Sunlit Needle & 238, 200, 111 & 阳光打在新生针叶上的暖金 \\
\bookcardswatch{222,166,32} & 蜂蜜松脂 & Honey Resin & 222, 166, 32 & 树皮裂口渗出的松脂金,是松树自愈的颜色 \\
\bookcardswatch{222,166,32} & 陈年金 & Aged Gold & 222, 166, 32 & 同色再现,使中调金更稳 \\
\bookcardswatch{177,120,133} & 幽影玫瑰 & Shadow Rose & 177, 120, 133 & 午后树荫里的玫瑰陶土,带紫调 \\
\bookcardswatch{167,88,26} & 树皮赭 & Bark Ochre & 167, 88, 26 & 树干主体的赭褐,是松树最诚实的颜色 \\
\bookcardswatch{43,47,34} & 林墨 & Forest Ink & 43, 47, 34 & 腐殖土与树根处的森林暗墨 \\
\bottomrule
\end{tabularx}\end{fullwidth}

\palettesubsec{松树年表的色彩映射}
\begin{palettebullets}
\item \textbf{Sunlit Needle} --- 阳针金是松树的「今天」——新生的、被光选中的、还没有故事的那一段时间
\item \textbf{Honey Resin} --- 蜂蜜松脂是松树的「伤口」——树脂封住裂痕,使损伤变成琥珀,是这组色彩最具时间感的一色
\item \textbf{Aged Gold} --- 陈年金与蜂蜜松脂共享色相,是中调金的「再说一遍」,把时间在颜色里复读一次
\item \textbf{Shadow Rose} --- 幽影玫瑰是松树投在地面的紫红色阴影,是树与土壤之间的呼吸界面
\item \textbf{Bark Ochre} --- 树皮赭是树干本体,是最不需要解释的色——它的赭褐就是松树这个生命体的物质本色
\item \textbf{Forest Ink} --- 林墨是松树的「根之色」,沉在腐殖土里,接近黑而仍属绿——它告诉我们,一棵活树的根,从不真正与黑相同
\end{palettebullets}
\par\bigskip
% ---------------------------------------------------------------------------
\bookpalettecardhead{6\enspace Provence Blue}
\palettesubsec{色彩哲学}
\noindent 普罗旺斯的蓝从不属于海——它属于薰衣草田尽头那片被夕阳压低饱和的蓝灰,属于午后石灰岩房屋窗框上经年褪色的孔雀蓝,属于一户人家晾在院里的旧亚麻床单与铁锈门把手共同呼吸的色温。这组配色拒绝地中海明信片式的浪漫,选择更内陆、更乡村、更被时间柔化过的法国南部蓝调——它的浪漫不在颜色的鲜艳里,而在颜色被使用过后留下的痕迹里。

\begin{fullwidth}\small\noindent
\begin{tabularx}{\linewidth}{@{}clllX@{}}
\toprule\rowcolor{booktableheadcolor}
\booktableheadcell{} & \booktableheadcell{中文名} & \booktableheadcell{原名} & \booktableheadcell{RGB} & \booktableheadcell{释义} \\
\midrule
\bookcardswatch{170,188,175} & 橄榄银灰 & Argent d'Olivier & 170, 188, 175 & 橄榄树叶背面的银灰,带绿调 \\
\bookcardswatch{137,156,154} & 雾镜蓝 & Miroir de Brume & 137, 156, 154 & 清晨雾气压在水面上的灰蓝 \\
\bookcardswatch{137,156,154} & 宿雨青 & Pluie Pass\'ee & 137, 156, 154 & 同色再现,巩固中调 \\
\bookcardswatch{110,124,139} & 石窗蓝 & Bleu de Volet & 110, 124, 139 & 石灰岩房屋窗框上的孔雀蓝褪色后的色 \\
\bookcardswatch{82,92,121} & 薰衣草夜 & Cr\'epuscule Lavande & 82, 92, 121 & 薰衣草田尽头被夕阳压低的暗紫蓝 \\
\bookcardswatch{53,66,94} & 橄榄夜空 & Nuit d'Olivier & 53, 66, 94 & 橄榄园上方夜空的深蓝,带土壤气息 \\
\bottomrule
\end{tabularx}\end{fullwidth}

\palettesubsec{南法乡村蓝调的色彩结构}
\begin{palettebullets}
\item \textbf{Olive Silver} --- 橄榄银灰是这组色彩里唯一的绿,是橄榄树叶背面在风里翻起的银光,是普罗旺斯空气的物质化
\item \textbf{Mist Mirror} --- 雾镜蓝是清晨水面的镜像,把蓝从天空借回地面,呼吸感最强的中调
\item \textbf{Spent Rain} --- 宿雨青与雾镜蓝同色,是「雨后第二日」的色温,刻意重复以加强配色的湿度
\item \textbf{Window Stone Blue} --- 石窗蓝把这组色彩从自然带回人造——它是被无数次刷过、又被无数次淡去的人手痕迹
\item \textbf{Lavender Dusk} --- 薰衣草夜是配色的灵魂——它把薰衣草从「紫」推向「暗紫蓝」,这一步,正是普罗旺斯之所以为普罗旺斯
\item \textbf{Olive Night} --- 橄榄夜空是这组色彩的根,把白日的所有蓝色收进了一片夜空,接近黑而仍是蓝
\end{palettebullets}
\par\bigskip
% ---------------------------------------------------------------------------
\bookpalettecardhead{7\enspace Fresco Blue}
\palettesubsec{色彩哲学}
\noindent 来自意大利文艺复兴时期教堂壁画的天空——那种被湿壁画(Fresco)技法吸入灰泥里、再被五百年烛烟氧化过的青蓝。它不是天空真实的颜色,而是匠人用青金石、孔雀石与石灰浆调出的「神圣的天空」,介于人间与天国之间。这组配色从最浅的天到最深的夜按梯度排列,模拟壁画从穹顶向地面跌落时蓝色的层层加深——它是一种被建筑收容过的天。

\begin{fullwidth}\small\noindent
\begin{tabularx}{\linewidth}{@{}clllX@{}}
\toprule\rowcolor{booktableheadcolor}
\booktableheadcell{} & \booktableheadcell{中文名} & \booktableheadcell{原名} & \booktableheadcell{RGB} & \booktableheadcell{释义} \\
\midrule
\bookcardswatch{166,224,244} & 穹顶浅天 & Cielo di Volta & 166, 224, 244 & 教堂穹顶最高处壁画的浅蓝 \\
\bookcardswatch{71,169,207} & 钴釉蓝 & Vetrina Cobalto & 71, 169, 207 & 青金石与石灰浆调和后的中调蓝 \\
\bookcardswatch{71,169,207} & 青金蓝 & Lapislazzuli & 71, 169, 207 & 同色再现,巩固壁画的中段 \\
\bookcardswatch{9,121,158} & 深殿蓝 & Azzurro d'Abside & 9, 121, 158 & 教堂半圆殿背景里五百年未褪的深蓝 \\
\bookcardswatch{4,75,102} & 夜祷蓝 & Azzurro della Veglia & 4, 75, 102 & 夜祷时灯火映在壁画下半部分的暗蓝 \\
\bookcardswatch{2,31,46} & 地宫蓝 & Azzurro della Cripta & 2, 31, 46 & 地宫与暗廊里壁画几近隐没的最深蓝 \\
\bottomrule
\end{tabularx}\end{fullwidth}

\palettesubsec{壁画天空的垂直色谱}
\begin{palettebullets}
\item \textbf{Vault Sky} --- 穹顶浅天是壁画最高处,光线打在灰泥上反射回来的浅蓝,是这组配色的「天」
\item \textbf{Cobalt Glaze} --- 钴釉蓝是中段的主调,是教堂内信众抬头时看见的那一抹饱满
\item \textbf{Lapis Renewed} --- 青金蓝与钴釉蓝同色,是「再次确认」——壁画的蓝从不依靠一次涂抹,它依靠层层叠加
\item \textbf{Apse Blue} --- 深殿蓝是教堂建筑里光线最难抵达的深处,壁画在那里凝固成不褪色的深蓝
\item \textbf{Vigil Blue} --- 夜祷蓝是夜晚的色温——同一面壁画,在油灯下显现出一种与白日截然不同的暗蓝
\item \textbf{Crypt Blue} --- 地宫蓝是壁画几乎被黑暗吞没的最深处,几近黑而仍属蓝——这是一组色彩中最沉默的「相信」
\end{palettebullets}
\par\bigskip
% ---------------------------------------------------------------------------
\bookpalettecardhead{8\enspace Monet}
\palettesubsec{色彩哲学}
\noindent 印象派的核心命题是:光本身有颜色。莫奈一生反复画同一座干草堆、同一座大教堂、同一池睡莲,只为捕捉「同一个对象在不同时刻拥有不同色彩」的事实。这组配色不复刻任何一张莫奈作品,而是从他全部画作里抽出最具印象派身份的六个色相:奶油画布、玫瑰干草、绿草地、夏池水、晨雾蓝、深苇影——它们之间的关系不是构图,而是同一束光在不同物质上的折射。

\begin{fullwidth}\small\noindent
\begin{tabularx}{\linewidth}{@{}clllX@{}}
\toprule\rowcolor{booktableheadcolor}
\booktableheadcell{} & \booktableheadcell{中文名} & \booktableheadcell{原名} & \booktableheadcell{RGB} & \booktableheadcell{释义} \\
\midrule
\bookcardswatch{247,244,213} & 画布奶油 & Toile Cr\`eme & 247, 244, 213 & 莫奈打底未上色画布的奶油底 \\
\bookcardswatch{211,150,140} & 玫瑰干草 & Foin Rose & 211, 150, 140 & 《干草堆》系列中黄昏时干草顶端的玫瑰 \\
\bookcardswatch{211,150,140} & 夕辉粉 & Lueur du Soir & 211, 150, 140 & 同色再现,延长黄昏的色温 \\
\bookcardswatch{131,153,88} & 草地绿 & Pr\'e Vert & 131, 153, 88 & 吉维尼花园午后草地的中调绿 \\
\bookcardswatch{16,86,102} & 池水蓝 & \'Etang Aigue & 16, 86, 102 & 《睡莲》系列水面下倒映的深蓝 \\
\bookcardswatch{10,51,35} & 芦苇暗影 & Ombre des Roseaux & 10, 51, 35 & 池边芦苇丛根部的深绿,接近墨 \\
\bottomrule
\end{tabularx}\end{fullwidth}

\palettesubsec{印象派光色的色彩结构}
\begin{palettebullets}
\item \textbf{Canvas Cream} --- 画布奶油是这组配色的「未完成」——是莫奈每一次起笔之前,画布本身的呼吸
\item \textbf{Hay Rose} --- 玫瑰干草是印象派最具识别性的一笔——它把「黄色」推向「玫瑰」,正是因为黄昏的光改写了一切固有色
\item \textbf{Evening Glow} --- 夕辉粉与玫瑰干草同色,是「延长黄昏」——印象派之所以反复画同一对象,正为多次捕捉这条色
\item \textbf{Meadow Green} --- 草地绿是配色的中调,是阳光打在新草上时的真实绿,既不冷也不暖
\item \textbf{Pond Aqua} --- 池水蓝是莫奈晚年《睡莲》中最具识别性的色相,把蓝从天空借给水,是这组配色的情绪转折点
\item \textbf{Reed Shadow} --- 芦苇暗影是这组配色的「根」,几近墨而仍是绿——印象派的暗部从不真正接受黑
\end{palettebullets}
\par\bigskip
% ---------------------------------------------------------------------------
\bookpalettecardhead{9\enspace Narcissus}
\palettesubsec{色彩哲学}
\noindent 不是希腊神话里的水仙,而是植物学意义上的水仙——一种从球根抽出花茎、在冬末春初开放的花。它的颜色谱系从苍白的鳞茎到金色的副冠、再到深褐的泥土根系,是一段「从地下到地面」的物质叙事。这组配色刻意去掉「绿叶」这一最显而易见的项,从而把视线集中在花本身从孕育、开放到回归大地的全过程,是一组以「时间纵轴」为结构的植物色谱。

\begin{fullwidth}\small\noindent
\begin{tabularx}{\linewidth}{@{}clllX@{}}
\toprule\rowcolor{booktableheadcolor}
\booktableheadcell{} & \booktableheadcell{中文名} & \booktableheadcell{原名} & \booktableheadcell{RGB} & \booktableheadcell{释义} \\
\midrule
\bookcardswatch{221,213,200} & 鳞茎乳 & Bulbus Lacteus & 221, 213, 200 & 刚切开的水仙球根截面的乳白 \\
\bookcardswatch{185,149,144} & 花瓣粉褐 & Petalum Roseobrunneum & 185, 149, 144 & 落瓣压在泥土里的灰玫瑰 \\
\bookcardswatch{185,149,144} & 退色玫瑰 & Rosa Defloruit & 185, 149, 144 & 同色再现,代表「再用一次」的语义 \\
\bookcardswatch{199,149,72} & 水仙金 & Aurum Narcissi & 199, 149, 72 & 副冠中央的暖金,是配色的标识色 \\
\bookcardswatch{190,108,26} & 蜂蜡铜 & Aes Mellitum & 190, 108, 26 & 球根外壳氧化后的橙褐,介于土与花之间 \\
\bookcardswatch{110,60,31} & 泥土深 & Humus Profundus & 110, 60, 31 & 水仙根系所在的腐殖土深褐,是根之色 \\
\bottomrule
\end{tabularx}\end{fullwidth}

\palettesubsec{水仙生命周期的色彩切片}
\begin{palettebullets}
\item \textbf{Bulb Milk} --- 鳞茎乳是水仙的「未开」,是植物在尚未成为花之前已经拥有的全部潜在颜色
\item \textbf{Petal Brown-Rose} --- 花瓣粉褐是花的「凋落」,花瓣离开花茎后被泥土与雨水共同改写的颜色
\item \textbf{Faded Rose} --- 退色玫瑰与花瓣粉褐同色,是「记忆中的同一片花瓣」——同一种凋落,在记忆里出现了第二次
\item \textbf{Narcissus Gold} --- 水仙金是这组配色的核心,是副冠中央那一抹只在春天初醒时出现的暖金,是「开放」本身的颜色
\item \textbf{Honey Bronze} --- 蜂蜡铜把花的金推向了土的赭,是花朵从开放走向凋落、从空中回到地下的过渡色
\item \textbf{Soil Deep} --- 泥土深是水仙的「来处」与「归处」,几近黑而保暖——这一组色彩告诉我们:任何盛开都从黑暗里来,也回去
\end{palettebullets}
\par\bigskip
% ---------------------------------------------------------------------------
\bookpalettecardhead{10\enspace Roman Empire}
\palettesubsec{色彩哲学}
\noindent 以罗马帝国的物质遗产为坐标:大理石、军团战袍、元老院托加、月桂冠、凯旋金、骨螺紫。六色不是设计而来,而是从罗马城本身长出来的——神庙的石、战场的血、广场的权、帝王的纹,共同构成一组配色的政治考古学。每一色都有可追溯的物质起源。

\begin{fullwidth}\small\noindent
\begin{tabularx}{\linewidth}{@{}clllX@{}}
\toprule\rowcolor{booktableheadcolor}
\booktableheadcell{} & \booktableheadcell{中文名} & \booktableheadcell{原名} & \booktableheadcell{RGB} & \booktableheadcell{释义} \\
\midrule
\bookcardswatch{236,232,225} & 大理石 & Marmor Carrariense & 236, 232, 225 & 卡拉拉大理石的冷白,略带暖灰,是神庙与雕像的物质 \\
\bookcardswatch{212,175,55} & 荣耀金 & Aurum Gloriae & 212, 175, 55 & 凯旋式黄金的中性金,介于铸币与神像镀金之间 \\
\bookcardswatch{74,110,65} & 桂冠 & Laurus Viridis & 74, 110, 65 & 月桂冠叶片的哑光深绿,取干燥叶片的沉稳色调 \\
\bookcardswatch{180,30,30} & 罗马军团 & Sagum Legionarium & 180, 30, 30 & 军团战袍(Sagum)的鲜烈战红,饱和而有力,象征铁与血 \\
\bookcardswatch{120,20,40} & 元老院 & Clavus Senatorius & 120, 20, 40 & 元老院托加袍缘的深沉酒红(Clavus),比军团红更内敛、更权贵 \\
\bookcardswatch{88,28,90} & 奥古斯都紫 & Purpura Tyria & 88, 28, 90 & 骨螺紫染料的历史色,价比黄金,专属皇权 \\
\bottomrule
\end{tabularx}\end{fullwidth}

\palettesubsec{设计思路}
\begin{palettebullets}
\item \textbf{Carrara Marble} --- 取自卡拉拉大理石的冷白底色,略带暖灰,呈现神庙与雕像的质感
\item \textbf{Legion Crimson} --- 军团战袍(Sagum)的鲜烈战红,饱和而有力,象征铁与血
\item \textbf{Senate Bordeaux} --- 元老院托加袍缘的深沉酒红(Clavus),比军团红更内敛、更权贵
\item \textbf{Laurel Viridis} --- 月桂冠叶片的哑光深绿,非翠绿,取干燥叶片的沉稳色调
\item \textbf{Gloria Aurum} --- 凯旋式黄金的中性金,介于铸币与神像镀金之间,不过暖也不过冷
\item \textbf{Tyrian Purple} --- 骨螺紫染料的历史色,价比黄金,专属皇权,深邃而神秘
\end{palettebullets}
\par\bigskip
% ---------------------------------------------------------------------------
\bookpalettecardhead{11\enspace Greece}
\palettesubsec{色彩哲学}
\noindent 以大理石白为底,民主蓝统领视觉重心,金色做高光点缀;陶器赤与橄榄绿构成一对冷暖平衡的中间调,葡萄紫作为最深的暗部锚点。大理石选用帕罗斯岛(Paros)产的暖白,比罗马卡拉拉更偏象牙,是帕特农神庙的实际用料;民主蓝取爱琴海深水区的中性蔚蓝,呼应雅典公民在广场(Agora)凝望的天与海。整体在地中海阳光下呈现出既庄重又充满生命力的古典张力。

\begin{fullwidth}\small\noindent
\begin{tabularx}{\linewidth}{@{}clllX@{}}
\toprule\rowcolor{booktableheadcolor}
\booktableheadcell{} & \booktableheadcell{中文名} & \booktableheadcell{原名} & \booktableheadcell{RGB} & \booktableheadcell{释义} \\
\midrule
\bookcardswatch{245,240,228} & 帕罗斯大理石 & \foreignname{Παριανός Μάρμαρος} \textperiodcentered\ Parianos Marmaros & 245, 240, 228 & 帕特农神庙的实际用料,比卡拉拉更偏象牙的暖白 \\
\bookcardswatch{212,175,55} & 荣耀金 & \foreignname{Χρυσός} \textperiodcentered\ Chrysos & 212, 175, 55 & 与罗马同源的中性金,呼应神明与宗教仪式的高光 \\
\bookcardswatch{48,105,175} & 民主蓝 & \foreignname{Κυανός Αγοράς} \textperiodcentered\ Kyanos Agoras & 48, 105, 175 & 爱琴海深水区的中性蔚蓝,雅典公民在广场凝望的天与海 \\
\bookcardswatch{188,82,38} & 陶器赤 & \foreignname{Αττικός Πηλός} \textperiodcentered\ Attikos Pelos & 188, 82, 38 & 红绘式陶器的底色,希腊人记录神话与英雄史诗最具标志性的媒介色 \\
\bookcardswatch{98,128,48} & 橄榄银绿 & \foreignname{Ελιά Αθηνάς} \textperiodcentered\ Elia Athenas & 98, 128, 48 & 雅典娜赠予城邦的橄榄树,和平与智慧的颜色 \\
\bookcardswatch{90,42,92} & 狄俄尼索斯葡萄紫 & \foreignname{Διονυσιακή Σταφυλή} \textperiodcentered\ Dionysiake Staphyle & 90, 42, 92 & 酿酒葡萄的深邃果紫,酒神、戏剧与狂欢节的灵魂色 \\
\bottomrule
\end{tabularx}\end{fullwidth}

\palettesubsec{配色逻辑}
\begin{palettebullets}
\item \textbf{Parian Marble} --- 大理石白作为底色,是希腊建筑与雕塑的物质本相——比罗马卡拉拉更暖、更象牙,带着地中海阳光过滤后的体温
\item \textbf{Agora Kyanos} --- 民主蓝统领整组色彩的视觉重心,既是爱琴海的天与海,也是雅典公民在广场上看见的「公共空间」的颜色
\item \textbf{Gloria Aurum} --- 荣耀金作为点缀,以最小面积承担神圣与神明的视觉重量,与罗马同源却用法不同——希腊的金更克制
\item \textbf{Attic Terracotta} --- 陶器赤是希腊的「叙事色」——红绘陶器把神话与英雄史诗压在这一色里,是希腊文化最具识别性的媒介
\item \textbf{Athena's Olive} --- 橄榄银绿与陶器赤构成冷暖平衡的中间调,是「智慧」与「叙事」的色彩对位
\item \textbf{Dionysian Grape} --- 葡萄紫是整组配色最深的暗部锚点,是酒神与戏剧的颜色——希腊文化的另一极:神圣之下的狂欢
\end{palettebullets}
\par\bigskip
% ---------------------------------------------------------------------------
\bookpalettecardhead{12\enspace Kanagawa}
\palettesubsec{色彩哲学}
\noindent 《神奈川冲浪里》的革命性在于葛饰北斎将当时刚传入日本的普鲁士蓝(ベロ藍)引入浮世绘,以单一色系的深浅变奏构建出整幅画面的戏剧张力,配合墨线轮廓与富士山的远景留白,形成极具压缩感的视觉美学。这组配色以普鲁士蓝家族为核心,从墨黑$\to$深海蓝$\to$浪蓝$\to$富士霞$\to$暮天$\to$波白构成一条完整的明度递进链,如同浪从深处涌起、在浪尖碎裂成雪沫的瞬间。

\begin{fullwidth}\small\noindent
\begin{tabularx}{\linewidth}{@{}clllX@{}}
\toprule\rowcolor{booktableheadcolor}
\booktableheadcell{} & \booktableheadcell{中文名} & \booktableheadcell{原名} & \booktableheadcell{RGB} & \booktableheadcell{释义} \\
\midrule
\bookcardswatch{237,233,222} & 波白 & \foreignname{波白} \textperiodcentered\ Nami-shiro & 237, 233, 222 & 浪尖破碎的泡沫,非纯白,带和纸的微暖底色 \\
\bookcardswatch{208,224,238} & 暮天 & \foreignname{暮天} \textperiodcentered\ Boten & 208, 224, 238 & 海平线处的天空渐变色,普鲁士蓝淡化至苍茫 \\
\bookcardswatch{150,186,210} & 富士霞 & \foreignname{富士霞} \textperiodcentered\ Fuji-gasumi & 150, 186, 210 & 远景富士山的淡蓝轮廓,虚化而沉静 \\
\bookcardswatch{26,78,132} & 普鲁士蓝 & \foreignname{ベロ藍} \textperiodcentered\ Bero-ai & 26, 78, 132 & 画面主体——普鲁士蓝浪身,江户最具标志性的舶来颜料 \\
\bookcardswatch{13,38,76} & 深海 & \foreignname{深海} \textperiodcentered\ Shinkai & 13, 38, 76 & 浪底最深处的墨蓝,近乎黑而仍保有蓝的冷意 \\
\bookcardswatch{29,25,35} & 墨 & \foreignname{墨} \textperiodcentered\ Sumi & 29, 25, 35 & 木版印刷的墨线,统摄全图骨架,非纯黑而含深紫底调 \\
\bottomrule
\end{tabularx}\end{fullwidth}

\palettesubsec{配色逻辑}
\begin{palettebullets}
\item \textbf{Nami-shiro} --- 波白是浪尖那一瞬间的「破碎」,带和纸的微暖底色,是整组色彩唯一的高光
\item \textbf{Boten} --- 暮天是海平线处天空的过渡色,把普鲁士蓝淡化至几乎消失,是画面「呼吸」的色
\item \textbf{Fuji-gasumi} --- 富士霞是远景富士山的淡蓝,是整幅画唯一的「退隐」,把焦点推回浪本身
\item \textbf{Bero-ai} --- 普鲁士蓝是浪身的主色,是这幅画的革命所在——一种单一外来颜料,撑起了整幅作品的戏剧张力
\item \textbf{Shinkai} --- 深海是浪底最深处的墨蓝,是「压迫感」与「深渊」的色彩载体,近乎黑而仍是蓝
\item \textbf{Sumi} --- 墨是木版印刷的骨架,统摄全图,是这组配色的「形」,没有它,六色蓝无可附着
\end{palettebullets}
\par\bigskip
% ---------------------------------------------------------------------------
\bookpalettecardhead{13\enspace Starry Night}
\palettesubsec{色彩哲学}
\noindent 梵高的《星月夜》并非印象派而是后印象派的巅峰之作——他不再捕捉光的瞬间,而是将内心的情绪直接燃烧进笔触。画面以钴蓝与群青为骨,用漩涡状的笔触赋予夜空以生命,月光与星芒的铬黄在蓝色的重压下迸裂成光爆,柏树如黑色火焰直刺苍穹,村庄的微光则是整幅画中唯一属于人间的暖。三层蓝构成一条由深至浅的情绪纵深轴,模拟梵高螺旋笔触的叠压节奏。

\begin{fullwidth}\small\noindent
\begin{tabularx}{\linewidth}{@{}clllX@{}}
\toprule\rowcolor{booktableheadcolor}
\booktableheadcell{} & \booktableheadcell{中文名} & \booktableheadcell{原名} & \booktableheadcell{RGB} & \booktableheadcell{释义} \\
\midrule
\bookcardswatch{240,208,68} & 月华铬黄 & Lumi\`ere Lunaire & 240, 208, 68 & 月亮与星芒的爆裂光晕,梵高常用铬黄表达极致的生命张力 \\
\bookcardswatch{198,140,52} & 村灯琥珀 & Lueurs du Village & 198, 140, 52 & 村庄窗口透出的暖黄,与冰冷蓝色星空形成全画唯一的冷暖对峙 \\
\bookcardswatch{105,155,200} & 晨曦苍蓝 & Aube Glac\'ee & 105, 155, 200 & 漩涡边缘的淡蓝过渡,深夜与月光交界处最微妙的呼吸感 \\
\bookcardswatch{48,96,165} & 星涡群青 & Tourbillon Outremer & 48, 96, 165 & 漩涡主体的群青蓝,后印象派用色的标志——非自然色而是情绪色 \\
\bookcardswatch{22,50,30} & 柏影墨绿 & Cypr\`es Nocturne & 22, 50, 30 & 前景柏树的深邃暗绿,几近于黑,是画面中唯一「拒绝光」的存在 \\
\bookcardswatch{20,36,88} & 深夜钴蓝 & Minuit Cobalt & 20, 36, 88 & 夜空最深处的底色,饱含压迫性的情绪重量,梵高精神世界的深渊底色 \\
\bottomrule
\end{tabularx}\end{fullwidth}

\palettesubsec{配色逻辑}
\begin{palettebullets}
\item \textbf{Lumi\`ere Lunaire} --- 月华铬黄是画面的「爆点」,月亮与星芒在蓝色重压下迸裂的光爆,是生命张力的极致
\item \textbf{Lueurs du Village} --- 村灯琥珀是画中唯一属于人间的暖,与冰冷的蓝色星空构成全画唯一的冷暖对峙
\item \textbf{Aube Glac\'ee} --- 晨曦苍蓝是过渡色,是深夜与月光交界处最微妙的呼吸感,介于压迫与解脱之间
\item \textbf{Tourbillon Outremer} --- 星涡群青是漩涡主体,是后印象派用色的标志——非自然色而是情绪色,是这组配色的「主旋」
\item \textbf{Cypr\`es Nocturne} --- 柏影墨绿是画面唯一「拒绝光」的存在,几近于黑,是视觉锚点,以最暗的色撑起整幅画的对比
\item \textbf{Minuit Cobalt} --- 深夜钴蓝是夜空最深处的情绪底色,是梵高精神世界的深渊,所有蓝色由它生长而来
\end{palettebullets}
\par\bigskip
% ---------------------------------------------------------------------------
\bookpalettecardhead{14\enspace A Thousand Li}
\palettesubsec{色彩哲学}
\noindent 《千里江山图》由北宋王希孟以矿物颜料层层积染而成——石青、石绿研磨后反复叠压,非一次上色,而是以「薄中见厚」的积染工艺令色泽沉入绢底,形成一种由内向外透光的矿物质感。六色之间非独立存在,而是共同构成一套山水呼吸的色彩生态。

\begin{fullwidth}\small\noindent
\begin{tabularx}{\linewidth}{@{}clllX@{}}
\toprule\rowcolor{booktableheadcolor}
\booktableheadcell{} & \booktableheadcell{中文名} & \booktableheadcell{原名} & \booktableheadcell{RGB} & \booktableheadcell{释义} \\
\midrule
\bookcardswatch{218,203,170} & 宋绢暖 & \foreignname{宋绢} \textperiodcentered\ S\`ong-ju\=an & 218, 203, 170 & 千年绢素氧化后的暖米底色,非颜料而是时间本身的色彩 \\
\bookcardswatch{110,165,195} & 烟渚色 & \foreignname{烟渚} \textperiodcentered\ Y\=an-zh\v{u} & 110, 165, 195 & 石青与铅白调和后的淡化层,用于远山轮廓与江面折光 \\
\bookcardswatch{183,118,45} & 砂岩赭 & \foreignname{砂赭} \textperiodcentered\ Sh\=a-zh\v{e} & 183, 118, 45 & 赭石(褐铁矿)与泥金混调,用于勾勒礁石肌理、屋舍廊桥与水波金线 \\
\bookcardswatch{96,148,110} & 岚霭绿 & \foreignname{岚霭绿} \textperiodcentered\ L\'an-\v{a}i-l\"u & 96, 148, 110 & 石绿渐次减量后的过渡色,表现中景植被与山腰云雾交融处 \\
\bookcardswatch{48,105,76} & 苍岭绿 & \foreignname{苍岭绿} \textperiodcentered\ C\=ang-l\v{i}ng-l\"u & 48, 105, 76 & 天然孔雀石的正色,层叠于山峦受光面,画面最具生命感的色调 \\
\bookcardswatch{35,78,112} & 霁山青 & \foreignname{霁山青} \textperiodcentered\ J\`i-sh\=an-q\=ing & 35, 78, 112 & 天然蓝铜矿研细后的深沉矿蓝,画中山体背阴面与水深处的主色 \\
\bottomrule
\end{tabularx}\end{fullwidth}

\palettesubsec{五大特征的色彩映射}
\begin{palettebullets}
\item \textbf{青绿主调} --- 霁山青与苍岭绿共同锚定全组色彩的主旋,一冷一暖,一阴一阳,相互制衡而非竞争
\item \textbf{层次渐变} --- 从霁山青$\to$烟渚色、苍岭绿$\to$岚霭绿,各形成一条明度递进链,模拟积染工艺的「薄中见厚」
\item \textbf{金赭点缀} --- 砂岩赭以最小的面积承担最大的温度调节职责,令青绿主调不至于流于清冷
\item \textbf{空灵雅致} --- 烟渚色与岚霭绿的介入,在浓重矿色之间打开呼吸的空间,令整组色不堵不闷
\item \textbf{古绢质感} --- 宋绢暖作为底色托起五色,如同千年绢素将所有浓烈都沉淀为一种悠远的克制
\end{palettebullets}
\par\bigskip
% ---------------------------------------------------------------------------
\bookpalettecardhead{15\enspace And Quiet Flows the Don}
\palettesubsec{色彩哲学}
\noindent 肖霍洛夫用二十年写完这部史诗,顿河哥萨克的命运从不在鲜亮之处——它埋在冻土里、锈在马刀上、凝在战壕边的黑血里、散在大雪漫过的苇荡中。这组配色拒绝一切「颜色」,只留下时间磨损之后的残余:不是蓝,是浊蓝;不是红,是暗血;不是黄,是枯草;不是灰,是铅灰。

\begin{fullwidth}\small\noindent
\begin{tabularx}{\linewidth}{@{}clllX@{}}
\toprule\rowcolor{booktableheadcolor}
\booktableheadcell{} & \booktableheadcell{中文名} & \booktableheadcell{原名} & \booktableheadcell{RGB} & \booktableheadcell{释义} \\
\midrule
\bookcardswatch{142,120,70} & 枯苇秋黄 & \foreignname{Полынь} \textperiodcentered\ Polyn & 142, 120, 70 & 顿河边被霜打过的苇草与苦艾,褪色的黄绿 \\
\bookcardswatch{110,108,104} & 草原铅灰 & \foreignname{Степной Пепел} \textperiodcentered\ Stepnoy Pepel & 110, 108, 104 & 冬日荒原的天与地混为一色,哥萨克男人眼神里的那种灰 \\
\bookcardswatch{118,86,52} & 黄土大地 & \foreignname{Донская Земля} \textperiodcentered\ Donskaya Zemlya & 118, 86, 52 & 顿河流域特有的栗钙土,带腐殖质气息的褐色 \\
\bookcardswatch{122,66,36} & 铁锈赭红 & \foreignname{Ржавое Железо} \textperiodcentered\ Rzhavoye Zhelezo & 122, 66, 36 & 哥萨克军刀、马镫、步枪枪托上的锈蚀,战争腐蚀金属的颜色 \\
\bookcardswatch{52,70,90} & 顿河浊蓝 & \foreignname{Мутный Дон} \textperiodcentered\ Mutnyy Don & 52, 70, 90 & 顿河从不清澈——裹挟泥沙与战死者的血,全书的时间底色 \\
\bookcardswatch{88,26,24} & 暗血战痕 & \foreignname{Запёкшаяся Кровь} \textperiodcentered\ Zapyokshayasya Krov & 88, 26, 24 & 凝固三日后渗入军衣的深褐红,血色命运在这里沉默 \\
\bottomrule
\end{tabularx}\end{fullwidth}

\palettesubsec{色彩结构的悲剧逻辑}
\begin{palettebullets}
\item \textbf{冷灰压顶} --- 草原铅灰与顿河浊蓝构成整组色彩的天空与水面,将所有暖色死死压在下方,如同历史对个体命运的碾压从未松手
\item \textbf{土褐承重} --- 黄土大地与枯苇秋黄是这片土地仅剩的「生」意,但它们同样去饱和、去鲜亮,美也是疲惫的美
\item \textbf{暗血锚底} --- 暗血战痕是全组最深的极色,不做主角,只在最沉的阴影处出现——如同《静静的顿河》里的死亡,从不戏剧化
\item \textbf{锈感贯穿} --- 铁锈赭红横跨土褐与暗血之间,是这组配色的「时间感」载体,它告诉观者:这里所有的一切,都已经过去很久了
\end{palettebullets}
\par\bigskip
% ---------------------------------------------------------------------------
\bookpalettecardhead{16\enspace Cyberpunk Edgerunners}
\palettesubsec{色彩哲学}
\noindent 《边缘行者》的视觉语言本质是过载美学——色彩不是装饰,是症状。大卫的精神崩溃伴随着荧光色彩的失控侵蚀,夜之城的霓虹不是繁荣的象征而是麻醉剂,Lucy 的蓝是这个世界里唯一真实的出口。这组配色的核心矛盾是:极暗的底色与极亮的荧光并置,产生几近生理不适的振动感。

\begin{fullwidth}\small\noindent
\begin{tabularx}{\linewidth}{@{}clllX@{}}
\toprule\rowcolor{booktableheadcolor}
\booktableheadcell{} & \booktableheadcell{中文名} & \booktableheadcell{原名} & \booktableheadcell{RGB} & \booktableheadcell{释义} \\
\midrule
\bookcardswatch{225,255,8} & 迷幻荧黄 & \foreignname{サイコ黄} \textperiodcentered\ Saiko-ki & 225, 255, 8 & 大卫夹克的灵魂色,赛博精神病发作时视觉过载之色 \\
\bookcardswatch{10,238,100} & 电路荧绿 & \foreignname{電脳緑} \textperiodcentered\ Denn\=o{}-ryoku & 10, 238, 100 & 网络骇入的代码流与改造体的皮下光纹 \\
\bookcardswatch{238,18,120} & 霓虹品红 & \foreignname{ネオン桃} \textperiodcentered\ Neon-momo & 238, 18, 120 & 夜之城最泛滥的霓虹,繁荣谎言的颜色 \\
\bookcardswatch{205,20,35} & 鲜血赤红 & \foreignname{エッジ赤} \textperiodcentered\ Edge-aka & 205, 20, 35 & 边缘行者的代价,街头枪战后水泥地上的颜色 \\
\bookcardswatch{15,32,98} & 露娜深蓝 & \foreignname{月の蒼} \textperiodcentered\ Tsuki-no-ao & 15, 32, 98 & Lucy 仰望月球的眼神,唯一不属于夜之城的颜色 \\
\bookcardswatch{14,8,28} & 夜都虚空 & \foreignname{夜都虚空} \textperiodcentered\ Yato Kok\=u{} & 14, 8, 28 & 夜之城街道底层的绝对黑暗,含深紫底调 \\
\bottomrule
\end{tabularx}\end{fullwidth}

\palettesubsec{色彩结构的叙事逻辑}
\begin{palettebullets}
\item \textbf{虚空压底} --- 夜都虚空作为绝对底色,饱和度趋近于零,是让其余五色得以「爆炸」的必要前提。没有黑暗,霓虹只是颜料
\item \textbf{冷暖裂变} --- 露娜深蓝与霓虹品红、迷幻荧黄构成最极端的冷暖对立:蓝是 Lucy 的月球梦,粉与黄是夜之城的谎言与疯狂
\item \textbf{荧光过载} --- 迷幻荧黄与电路荧绿是赛博精神病美学的视觉核心,两色并置时产生几近生理不适的振动感
\item \textbf{红色终章} --- 鲜血赤红在色板中不做主角,却是最诚实的颜色——像结局一样,无可回避
\end{palettebullets}
\par\bigskip
% ---------------------------------------------------------------------------
\bookpalettecardhead{17\enspace The Grand Budapest Hotel}
\palettesubsec{色彩哲学}
\noindent 韦斯·安德森的配色从不是写实的,而是经过精密校准的情感滤镜——《布达佩斯大饭店》的粉色不是粉色,是一整个正在消逝的旧欧洲贵族世界被压缩成的颜色;紫色不是紫色,是古斯塔夫先生永不失礼的最后尊严。每一种颜色都经过刻意的去真实化处理。

\begin{fullwidth}\small\noindent
\begin{tabularx}{\linewidth}{@{}clllX@{}}
\toprule\rowcolor{booktableheadcolor}
\booktableheadcell{} & \booktableheadcell{中文名} & \booktableheadcell{原名} & \booktableheadcell{RGB} & \booktableheadcell{释义} \\
\midrule
\bookcardswatch{247,235,215} & 香草奶油 & Cr\`eme Vanille & 247, 235, 215 & 饭店内壁、Mendl 蛋糕胚、信笺底色,童话世界的空气本身 \\
\bookcardswatch{168,178,196} & 阿尔卑斯雾蓝 & Brume Alpine & 168, 178, 196 & 积雪山景、冬日车厢窗外的低饱和冷调,情绪降温剂 \\
\bookcardswatch{237,148,158} & 美酒蔷薇 & Rose M\'endl & 237, 148, 158 & 饭店外墙与 Mendl 甜品盒的标志粉,掺了一丝尘埃的旧玫瑰色 \\
\bookcardswatch{176,136,60} & 古铜暗金 & Dor\'e Antique & 176, 136, 60 & 镀金镜框、门把手、制服肩章的年岁感,旧贵族的光泽 \\
\bookcardswatch{128,88,148} & 执行者紫 & Violet Concierge & 128, 88, 148 & 古斯塔夫先生制服的紫,兼具权威与荒诞 \\
\bookcardswatch{143,52,65} & 复古酒绛 & Bordeaux Vintage & 143, 52, 65 & 深色场景的戏剧性锚点——图书馆皮革、红酒、密谋与危险 \\
\bottomrule
\end{tabularx}\end{fullwidth}

\palettesubsec{五重美学法则的色彩映射}
\begin{palettebullets}
\item \textbf{马卡龙柔色} --- 蔷薇粉与执行者紫是主视觉担当,高饱和但经过刻意的明度压制,像马卡龙饼壳——色彩分明却绝不刺眼
\item \textbf{奶油暖底} --- 香草奶油以最大留白统摄全局,是韦斯·安德森式对称构图里的「负空间」,令其他五色不至于拥挤
\item \textbf{复古暗金} --- 古铜暗金专门负责「贵族感」的重量,它不参与童话叙事,只在边框与细节处发光
\item \textbf{冷调克制} --- 阿尔卑斯雾蓝是全组唯一的冷色,低饱和度让它成为调节器,在甜腻的暖色系中打开一扇透气的窗
\item \textbf{甜而不腻} --- 复古酒绛是整组配色最精妙的设计:一粒深沉的暗色锚点,像甜点里的一滴苦杏仁酒,让所有的轻盈都有了重量
\end{palettebullets}
\par\bigskip
% ---------------------------------------------------------------------------
\bookpalettecardhead{18\enspace Renaissance Florence}
\palettesubsec{色彩哲学}
\noindent 佛罗伦萨的色彩并非设计而来,而是从物质本身生长出来的——青金石研磨成颜料、铅白调入蛋黄成坦培拉底料、金箔贴于木质祭坛、陶土在托斯卡纳窑火中烧成穹顶砖瓦。这六种颜色皆有真实的物质出处:它们曾是颜料、矿石、建材、木料,是波提切利、乔托与多纳泰罗共同呼吸的色谱。

\begin{fullwidth}\small\noindent
\begin{tabularx}{\linewidth}{@{}clllX@{}}
\toprule\rowcolor{booktableheadcolor}
\booktableheadcell{} & \booktableheadcell{中文名} & \booktableheadcell{原名} & \booktableheadcell{RGB} & \booktableheadcell{释义} \\
\midrule
\bookcardswatch{238,228,208} & 翡冷翠象牙 & Avorio Fiorentino & 238, 228, 208 & 卡拉拉大理石截面与坦培拉画板的打底铅白,含蛋黄与亚麻油的暖意 \\
\bookcardswatch{190,148,52} & 祭坛暖金 & Oro dell'Altare & 190, 148, 52 & 圣母像背光与美第奇金饰的金箔色,五百年空气氧化后内敛发光的旧金 \\
\bookcardswatch{172,84,50} & 穹顶陶红 & Cotto Brunellesco & 172, 84, 50 & 布鲁内莱斯基穹顶砖瓦与佛罗伦萨屋檐的标志陶土红 \\
\bookcardswatch{54,80,60} & 柏影墨绿 & Verde Cipresso & 54, 80, 60 & 托斯卡纳丘陵柏树与铜绿颜料(Verdigris)的混调 \\
\bookcardswatch{46,70,118} & 青金矿蓝 & Oltremare di Lapislazzuli & 46, 70, 118 & 产自阿富汗矿山的青金石研磨所得,价比黄金,专用于圣母袍服 \\
\bookcardswatch{85,54,34} & 胡桃深褐 & Noce Toscano & 85, 54, 34 & 佛罗伦萨工坊核桃木画板与深色底釉的颜色 \\
\bottomrule
\end{tabularx}\end{fullwidth}

\palettesubsec{四重气质的色彩映射}
\begin{palettebullets}
\item \textbf{神圣庄重} --- 青金矿蓝承担全组的精神重量。青金石在文艺复兴时期价比黄金,非普通委托可用,它的存在本身即是神圣的声明;祭坛暖金以最克制的方式散发光晕,二者共同构建宗教艺术的仪式感
\item \textbf{温润华贵} --- 翡冷翠象牙与祭坛暖金是贵族底色的两面——一个是触感温润的大理石与绢布,一个是美第奇家族以财富供养艺术的金属记忆;陶土红以人间温度调和神圣的冷峻
\item \textbf{古典醇厚} --- 六色均经过刻意的去现代化处理:无一纯色,皆含矿物杂质的「不纯粹」——这正是天然颜料与合成色料最根本的区别。厚重来自饱和度的克制,醇厚来自每个颜色内部的复杂性
\item \textbf{人文雅致} --- 柏影墨绿与胡桃深褐是人文主义的色彩锚点:绿来自托斯卡纳的自然山丘,褐来自工坊的核桃木画板,它们将神圣从天空拉回大地
\end{palettebullets}
\par\bigskip
% ---------------------------------------------------------------------------
\bookpalettecardhead{19\enspace Soviet Avant-Garde}
\palettesubsec{色彩哲学}
\noindent 苏联先锋派从不「设计配色」——他们宣告颜色。罗钦科说色彩必须有任务,李西茨基说形式即政治,马列维奇说至上主义的黑色正方形是零度的绝对。这组配色拒绝调和、拒绝渐变、拒绝装饰性的过渡,每种颜色都是一个硬边宣言:红色不解释自己,黑色不道歉,灰色不讨好任何人。

\begin{fullwidth}\small\noindent
\begin{tabularx}{\linewidth}{@{}clllX@{}}
\toprule\rowcolor{booktableheadcolor}
\booktableheadcell{} & \booktableheadcell{中文名} & \booktableheadcell{原名} & \booktableheadcell{RGB} & \booktableheadcell{释义} \\
\midrule
\bookcardswatch{225,218,205} & 粗纸白 & \foreignname{ГАЗЕТА} \textperiodcentered\ Gazeta & 225, 218, 205 & 《真理报》与构成主义宣传册的新闻纸白,含油墨渗透后的暖灰底调 \\
\bookcardswatch{200,150,32} & 宣传赭黄 & \foreignname{ПЛАКАТ} \textperiodcentered\ Plakat & 200, 150, 32 & 克鲁齐斯摄影蒙太奇海报与农业集体化宣传画的赭黄 \\
\bookcardswatch{118,114,108} & 混凝土灰 & \foreignname{БЕТОН} \textperiodcentered\ Beton & 118, 114, 108 & 塔特林塔、莫斯科工厂厂房、五年计划基建工地的浇筑混凝土色 \\
\bookcardswatch{50,80,138} & 机械蓝 & \foreignname{ЧЕРТЁЖ} \textperiodcentered\ Chertyozh & 50, 80, 138 & 工程蓝图与苏联工人阶级制服的冷蓝,去饱和处理 \\
\bookcardswatch{196,28,28} & 列宁红 & \foreignname{КРАСНЫЙ} \textperiodcentered\ Krasnyy & 196, 28, 28 & 布尔什维克旗帜与罗钦科海报的革命红,刻意压低饱和度令其沉重如铁 \\
\bookcardswatch{24,20,18} & 铸铁黑 & \foreignname{ЧЁРНЫЙ} \textperiodcentered\ Chyornyy & 24, 20, 18 & 马列维奇《黑色正方形》的绝对黑,虚无与确定性同时存在的黑 \\
\bottomrule
\end{tabularx}\end{fullwidth}

\palettesubsec{六重美学法则的色彩映射}
\begin{palettebullets}
\item \textbf{红黑强撞} --- 列宁红与铸铁黑是整组配色的结构核心,二者之间无过渡、无缓冲,直接对撞——对比本身即是论点
\item \textbf{工业灰的中轴} --- 混凝土灰不参与任何戏剧,它只是站在那里承重,如同五年计划中那些无名工人,是这组配色的脊柱而非焦点
\item \textbf{机械蓝的精确} --- 机械蓝刻意压低至近乎去色的状态,与列宁红形成冷暖的意识形态对立:红是革命的激情,蓝是建设的纪律
\item \textbf{粗纸白的民主性} --- 新闻纸白是这组配色里唯一的「反精英」色,它拒绝象牙白的贵族温润,选择廉价印刷品的粗糙质感
\item \textbf{赭黄的未竟性} --- 宣传赭黄是全组配色中最具悲剧性的颜色:它是海报上的丰收麦穗、是乌托邦宣传画里永远灿烂的明天,是一个承诺的颜色而非现实的颜色
\end{palettebullets}
\par\bigskip
% ---------------------------------------------------------------------------
\bookpalettecardhead{20\enspace Constantinople}
\palettesubsec{色彩哲学}
\noindent 君士坦丁堡是人类历史上唯一一座同时属于两个大陆、三种文明、四种宗教的城市——它的颜色从不单纯。圣索菲亚大教堂的穹顶金光里叠压着希腊工匠、罗马工程、东方神学;博斯普鲁斯海峡的深蓝里混着黑海的寒流与爱琴海的温暖;帝王紫是用地中海骨螺提炼出的奢华。

\begin{fullwidth}\small\noindent
\begin{tabularx}{\linewidth}{@{}clllX@{}}
\toprule\rowcolor{booktableheadcolor}
\booktableheadcell{} & \booktableheadcell{中文名} & \booktableheadcell{原名} & \booktableheadcell{RGB} & \booktableheadcell{释义} \\
\midrule
\bookcardswatch{236,225,207} & 普罗科尼索象牙 & \foreignname{Λευκός} \textperiodcentered\ Fildi\c{s}i & 236, 225, 207 & 马尔马拉海小岛出产的普罗科尼索大理石,铺满圣索菲亚内殿地面 \\
\bookcardswatch{195,150,42} & 圣索菲亚金 & \foreignname{Χρυσός} \textperiodcentered\ Alt\i n & 195, 150, 42 & 圣索菲亚穹顶马赛克金箔的光,一千五百年烛烟熏染后内敛发光的旧金 \\
\bookcardswatch{170,112,50} & 狄奥多西赭 & \foreignname{\'{Ω}χρα} \textperiodcentered\ Toprak & 170, 112, 50 & 狄奥多西城墙砖石与博斯普鲁斯岸边老建筑的赭土色 \\
\bookcardswatch{50,76,58} & 柏树暗绿 & \foreignname{Κυπαρίσσι} \textperiodcentered\ Selvi & 50, 76, 58 & 君士坦丁堡山丘柏树与拜占庭绿色斑岩柱础的深沉暗绿 \\
\bookcardswatch{102,32,65} & 帝王骨螺紫 & \foreignname{Πορφύρα} \textperiodcentered\ Mor & 102, 32, 65 & 拜占庭皇帝专属的骨螺紫,提炼一克需杀万只骨螺 \\
\bookcardswatch{30,52,88} & 博斯普鲁斯蓝 & \foreignname{Βόσπορος} \textperiodcentered\ Bo\u{g}az & 30, 52, 88 & 海峡傍晚收拢光线后的深沉蓝,黑海与爱琴海在此交汇 \\
\bottomrule
\end{tabularx}\end{fullwidth}

\palettesubsec{文明交汇的色彩结构}
\begin{palettebullets}
\item \textbf{东西冷暖轴} --- 博斯普鲁斯深蓝(西方理性、海洋秩序)与狄奥多西赭(东方大地、安纳托利亚土壤)构成配色的东西文明轴线;二者温度相反,却在君士坦丁堡这座城市里共存了千年而未曾消解
\item \textbf{神圣三角} --- 圣索菲亚金、帝王骨螺紫、普罗科尼索象牙共同构建宗教神圣感的色彩空间:金是神的光,紫是神在人间的代理人,象牙是神居住其中的建筑肌肤
\item \textbf{紫红的文明桥接} --- 帝王骨螺紫在整组配色中承担最复杂的角色:它同时属于罗马帝国的皇权传统、拜占庭基督教的神学色彩,以及波斯-东方宫廷的奢华美学
\item \textbf{暗绿的历史纵深} --- 柏树暗绿以最低调的姿态承载最漫长的时间跨度:它同时出现在拜占庭墓地柏树、奥斯曼庭园绿化与伊斯兰清真寺的铜绿穹顶上
\item \textbf{金辉内敛} --- 整组配色刻意将圣索菲亚金的饱和度压低至矿物感区间,令其发光而不炫目——这是拜占庭马赛克艺术的核心逻辑:金不是装饰,是光本身的物质化
\end{palettebullets}
\par\bigskip
% ---------------------------------------------------------------------------
\bookpalettecardhead{21\enspace France}
\palettesubsec{色彩哲学}
\noindent 法国的颜色从不是单纯的颜色——它是一种世界观的物质显现。共和红里藏着断头台与玫瑰,深蓝里住着笛卡尔的理性与莫奈的晨雾,香槟白是凡尔赛大理石与左岸咖啡馆信纸同一种底色。这组配色试图捕捉法国精神最核心的内在矛盾:极度理性与极度浪漫在同一块土地上共生。

\begin{fullwidth}\small\noindent
\begin{tabularx}{\linewidth}{@{}clllX@{}}
\toprule\rowcolor{booktableheadcolor}
\booktableheadcell{} & \booktableheadcell{中文名} & \booktableheadcell{原名} & \booktableheadcell{RGB} & \booktableheadcell{释义} \\
\midrule
\bookcardswatch{180,32,40} & 共和红 & Rouge Marianne & 180, 32, 40 & 三色旗之红,但非正红——波尔多红酒的深度与自由女神帽的激情的中间值 \\
\bookcardswatch{232,222,205} & 香槟象牙 & Ivoire Champagne & 232, 222, 205 & 非纯白——香槟气泡的暖调、卢浮宫石灰岩的温度、左岸咖啡馆信纸 \\
\bookcardswatch{44,74,138} & 共和深蓝 & Bleu R\'epublique & 44, 74, 138 & 三色旗之蓝,取塞夫勒瓷器皇家蓝的深度——启蒙理性的蓝 \\
\bookcardswatch{143,108,148} & 普罗旺斯薰衣草 & Lavande Proven\c{c}ale & 143, 108, 148 & 南法薰衣草田在午后阳光与灰白石灰土之间呈现的真实颜色 \\
\bookcardswatch{142,133,58} & 凡尔赛橄榄金 & Olive Dor\'ee & 142, 133, 58 & 凡尔赛宫正式花园修剪黄杨篱的金绿——古典法式园林美学 \\
\bookcardswatch{122,126,132} & 巴黎锌灰 & Zinc Parisien & 122, 126, 132 & 奥斯曼改造后巴黎屋顶锌板在阴天的反光色 \\
\bottomrule
\end{tabularx}\end{fullwidth}

\palettesubsec{法兰西精神的色彩结构}
\begin{palettebullets}
\item \textbf{三色旗的再解读} --- 共和红、香槟象牙、共和深蓝并非原版三色旗的复刻,而是对三色精神的去符号化再提炼:红被赋予波尔多的重量,白被赋予时间的温度,蓝被赋予启蒙的深度——旗帜是政治的,这组色是文化的
\item \textbf{南北气质的张力} --- 普罗旺斯薰衣草与巴黎锌灰构成法国南北最核心的精神对立:南法是感官、土地、阳光与慵懒的浪漫;巴黎是理性、几何、灰色天空与克制的优雅
\item \textbf{先锋撞色的潜伏} --- 凡尔赛橄榄金与共和深蓝的并置是全组最具先锋张力的组合:一个属于太阳王的古典秩序,一个属于共和革命的理性锋芒——色相接近却语义对立
\item \textbf{克制高级的实现路径} --- 六色均经过刻意的饱和度压制:红不燃烧,蓝不冰冷,紫不妖艳,金不张扬,灰不死板,白不刺目。法国人所理解的「高级」,从来不是色彩的强度,而是色彩恰到好处地收敛自身
\item \textbf{土地温暖的隐线} --- 香槟象牙与普罗旺斯薰衣草共同编织这组配色的「土地感」底线:前者是石灰岩与小麦的暖,后者是干燥薰衣草与灰白土壤的暖——法国的浪漫从不悬浮在空中
\end{palettebullets}
\par\bigskip
% ---------------------------------------------------------------------------
\bookpalettecardhead{22\enspace Kyoto}
\palettesubsec{色彩哲学}
\noindent 京都的颜色从不主动开口——它们等待被看见。禅院的苔绿不争,枯山水的白砂不语,町家杉木的茶褐在雨中沉默,暗绀的夜空压着千年古都不作声。唯有朱红鸟居,在一切克制之中以极小的面积燃烧最大的温度,像茶室里一枝单花,冷寂里的那一点热烈。

\begin{fullwidth}\small\noindent
\begin{tabularx}{\linewidth}{@{}clllX@{}}
\toprule\rowcolor{booktableheadcolor}
\booktableheadcell{} & \booktableheadcell{中文名} & \booktableheadcell{原名} & \booktableheadcell{RGB} & \booktableheadcell{释义} \\
\midrule
\bookcardswatch{234,226,212} & 绢白 & \foreignname{絹白} \textperiodcentered\ Kinushiro & 234, 226, 212 & 西阵织绢布与和纸漉纸的底色,蚕丝与楮树皮的天然暖调 \\
\bookcardswatch{110,84,64} & 枯茶褐 & \foreignname{枯茶} \textperiodcentered\ Karacha & 110, 84, 64 & 町家百年杉木柱在时间与炉烟中形成的褐色,千利休「侘び」的颜色 \\
\bookcardswatch{204,184,178} & 樱灰粉 & \foreignname{桜鼠} \textperiodcentered\ Sakura-nezumi & 204, 184, 178 & 落樱覆于石灯笼上被雨水晕开后的颜色,物哀的精确定义 \\
\bookcardswatch{116,126,120} & 青竹灰 & \foreignname{青竹鼠} \textperiodcentered\ Aotake-nezumi & 116, 126, 120 & 岚山竹林晨雾中的竹节色,冷暖平衡的禅理体现 \\
\bookcardswatch{80,94,68} & 苔色 & \foreignname{苔色} \textperiodcentered\ Koke-iro & 80, 94, 68 & 龙安寺枯山水石组底部苔藓的深绿,时间留下的暗哑印记 \\
\bookcardswatch{34,40,66} & 紺 & \foreignname{紺} \textperiodcentered\ Kon & 34, 40, 66 & 能剧服装与京友禅夜色图案的深靛蓝,冥想状态本身的颜色 \\
\bottomrule
\end{tabularx}\end{fullwidth}

\palettesubsec{京都精神的色彩结构}
\begin{palettebullets}
\item \textbf{留白为骨} --- 绢白与樱灰粉共同构建配色的呼吸空间,前者是空无、负形、茶室墙面,后者是将消未消的感知边界——京都美学不在于填满,而在于知道哪里应当什么都没有
\item \textbf{冷暖天平} --- 六色底中,枯茶褐与苔绿属暖调,暗绀与青竹灰属冷调,樱灰粉与绢白居中游移——四季平衡,禅院冷寉与町家温暖在此精确对位
\item \textbf{幽玄的低饱和逻辑} --- 所有色彩均经过刻意的去鲜亮处理,模拟京都特有的盆地晨雾对色彩的柔化效果——雾不消色,只让色彩向内收,这正是「幽玄」的物理机制
\item \textbf{水墨留白的东方结构} --- 整组配色可理解为一幅水墨:暗绀是浓墨,苔绿是中墨,枯茶褐是淡墨,青竹灰是宿墨,樱灰粉是墨的晕散,绢白是纸——朱红是画师最后蘸朱砂落下的那一点钤印
\end{palettebullets}
\par\bigskip
% ---------------------------------------------------------------------------
\bookpalettecardhead{23\enspace Siamese Dream}
\palettesubsec{色彩哲学}
\noindent 《Siamese Dream》的封面本身就是一张配色宣言——两个女孩在柔焦的暖白光晕里相依,背景的深绿若隐若现,橙黄的光线像记忆里某个下午,美好得不真实。这不是摇滚专辑封面惯用的高对比冲击美学,而是一种刻意制造的梦境质感:过曝、柔化、带着胶片颗粒的温柔失真——如同 Corgan 把最原始的情感录进了最失真的吉他音墙里。

\begin{fullwidth}\small\noindent
\begin{tabularx}{\linewidth}{@{}clllX@{}}
\toprule\rowcolor{booktableheadcolor}
\booktableheadcell{} & \booktableheadcell{中文名} & \booktableheadcell{原名} & \booktableheadcell{RGB} & \booktableheadcell{释义} \\
\midrule
\bookcardswatch{238,226,208} & 终章暖白 & Luna & 238, 226, 208 & 封面的暖白光晕,专辑最后一首歌的余韵——不是纯白,是胶片过曝后仍保留的一丝人间温度,是梦醒前最后的柔光 \\
\bookcardswatch{182,180,193} & 裸露灰紫 & Disarm & 182, 180, 193 & 洗白的灰紫,Disarm 是整张专辑里防御最低的时刻,弦乐与铃声剥去了所有失真——这是玻璃碎裂前的颜色,透明、脆、带着童年创伤的寒意 \\
\bookcardswatch{198,158,65} & 可疑旧金 & Today & 198, 158, 65 & 沾了尘的旧金色,Today 听起来是全碟最明媚的旋律,却是 Corgan 最黑暗时刻的产物——这抹金因此必须是褪色的、带着微苦的甜,美丽而不可信任 \\
\bookcardswatch{150,145,135} & 悬浮烟灰 & Hummer & 150, 145, 135 & 烟灰褐,Hummer 用层层叠叠的音墙制造一种失重的悬浮感——这是梦游者眼中的颜色,既不清醒也不沉睡,是 shoegaze 美学最诚实的物质形态 \\
\bookcardswatch{188,105,40} & 低温烧橙 & Mayonaise & 188, 105, 40 & 深而沉的烧橙,封面光源的颜色——Mayonaise 是全碟最慢、最痛、最私密的曲目,这抹橙不是热情而是积压已久的情感在低温中缓慢焦化的颜色 \\
\bookcardswatch{42,68,48} & 深渊暗绿 & Soma & 42, 68, 48 & 封面背景若隐若现的深绿,九分钟的史诗长曲从轻柔攀爬至崩塌——这是有机物在黑暗中生长的颜色,苔藓与腐叶,美丽与窒息共生 \\
\bottomrule
\end{tabularx}\end{fullwidth}

\palettesubsec{青春期情感的悖论结构}
\begin{palettebullets}
\item \textbf{梦境包裹与失真美学} --- Luna 的暖白与 Hummer 的烟灰共同构成专辑的“失焦层”——一切痛苦都先被过曝、被柔化、被音墙包住,然后才允许被听见。这不是回避痛苦,而是让痛苦以可承受的密度释放
\item \textbf{美丽容器装最重的内容} --- Today 的旧金与 Disarm 的灰紫是悖论的两极:前者用明媚旋律包裹自毁倾向,后者用童谣式的铃声讲述童年创伤——Corgan 始终用最美的容器装最重的内容,这是 Siamese Dream 最核心的情感策略
\item \textbf{情感核心的低温焦化} --- Mayonaise 的烧橙是整组配色的情感核心(pal5=primary 也正落在此处)——不是燃烧的红,而是已经燃烧过、被时间慢烤、被自身重量压低的橙。它告诉读者:这里的热度不来自爆发,而来自长期不灭
\item \textbf{深渊作为最深锚点} --- Soma 的暗绿(pal6=accent)是六色中唯一的“黑色替身”——它不是纯黑,而是有机物的黑,是苔藓、腐叶、深林。整组配色因此从不依赖纯黑来收尾,而是用一种“仍在生长的黑暗”把整张专辑稳稳压住
\item \textbf{回避高饱和与强对比的渗透美学} --- 六色全部退避鲜亮,因为 Siamese Dream 的情绪从不“爆炸”——它“渗透”。没有任何一色试图夺取注意力,反而每一色都在以最低的姿态长时间地停留在视觉里,如同 shoegaze 的吉他延音
\end{palettebullets}
